\documentclass[12pt]{article}
\usepackage[utf8]{inputenc}
\usepackage[T1]{fontenc}
\usepackage[margin=1in]{geometry}
\usepackage{setspace}
\usepackage{xcolor}
\usepackage[colorlinks=true, urlcolor=blue!60!black, linkcolor=blue!60!black]{hyperref}
\setlength{\parindent}{2em}
\setlength{\parskip}{1.25em}
\renewcommand{\baselinestretch}{1.0}
\title{\textit{South Dakota v. Dole}, 483 U.S. 203 (1987)}
\author{Case Brief}
\date{}
\begin{document}
\maketitle

% Case Brief: South Dakota v. Dole, 483 U.S. 203 (1987)
% Chapter: Federal Legislative Power
% Prepared for Seth J. Chandler's Constitutional Law Briefing Book

\section{Case Metadata}

\begin{itemize}
  \item \textbf{Case Name:} \href{https://supreme.justia.com/cases/federal/us/483/203/}{\textit{South Dakota v. Dole}} (Elizabeth Dole, Secretary of Transportation)
  \item \textbf{Citation:} 483 U.S. 203 (1987); Docket No. 86-260 (\href{https://tile.loc.gov/storage-services/service/ll/usrep/usrep483/usrep483203/usrep483203.pdf}{Official U.S. Reports PDF, Library of Congress})
  \item \textbf{Dates:} Argued April 28, 1987; Decided June 23, 1987
  \item \textbf{Court Below:} United States Court of Appeals for the Eighth Circuit, \href{https://caselaw.findlaw.com/court/us-8th-circuit/114759006.html}{791 F.2d 628 (8th Cir. 1986)} (affirming the United States District Court for the District of South Dakota)
  \item \textbf{Disposition:} Affirmed (7--2)
  \item \textbf{Vote and Alignment:}
  \begin{itemize}
    \item \textbf{Majority (7):} Chief Justice William H. Rehnquist (author), joined by Justices Byron R. White, Thurgood Marshall, Harry A. Blackmun, Lewis F. Powell, Jr., John Paul Stevens, and Antonin Scalia.
    \item \textbf{Concurrences (0):} None.
    \item \textbf{Dissents (2):} Justice William J. Brennan, Jr. (separate dissenting opinion, 483 U.S. at 212); Justice Sandra Day O'Connor (dissenting opinion, 483 U.S. at 212--218).
  \end{itemize}
  \item \textbf{Constitutional Provisions at Issue:} Taxing and Spending Clause, U.S. Const. art. I, \S~8, cl. 1; Twenty-first Amendment, \S~2; Tenth Amendment (structural federalism principles).
  \item \textbf{Statutes at Issue:} National Minimum Drinking Age Act of 1984 (23 U.S.C. \S~158 (1982 ed., Supp. III)); S.D. Codified Laws \S~35-6-27 (1986).
  \item \textbf{Primary Doctrinal Topic:} Conditional Federal Spending Power; Limits on Congressional Inducements to States; Interplay Between the Spending Clause and Structural Constitutional Limits.
\end{itemize}

\section{Facts and Procedural History}

In 1984, Congress enacted 23 U.S.C. \S~158 (1982 ed., Supp. III), commonly known as the National Minimum Drinking Age Act (or National Minimum Drinking Age Amendment), responding to findings by the Presidential Commission on Drunk Driving (1983) that inconsistent state minimum drinking ages created ``an incentive to drink and drive'' by encouraging young adults to commute across state borders to purchase alcohol (\href{https://supreme.justia.com/cases/federal/us/483/203/}{483 U.S. at 209}). Rather than enacting a direct federal prohibition on underage alcohol sales, \S~158 directed the Secretary of Transportation to withhold 5\% of federal highway funds otherwise allocable under specified federal highway grant programs from any state ``in which the purchase or public possession . . . of any alcoholic beverage by a person who is less than twenty-one years of age is lawful'' (\href{https://supreme.justia.com/cases/federal/us/483/203/}{483 U.S. at 205, 211}).

At the time, South Dakota law permitted persons 19 years of age or older to purchase beer containing up to 3.2\% alcohol (S.D. Codified Laws \S~35-6-27 (1986)). Facing the loss of 5\% of its covered federal highway funds, South Dakota sued Secretary of Transportation Elizabeth Dole in the United States District Court for the District of South Dakota, seeking a declaratory judgment that \S~158 exceeded Congress's authority under the Spending Clause (U.S. Const. art. I, \S~8, cl. 1) and violated the Twenty-first Amendment's reservation of state authority over alcoholic beverages (\href{https://supreme.justia.com/cases/federal/us/483/203/}{483 U.S. at 205}).

The District Court rejected South Dakota's constitutional challenges, and the United States Court of Appeals for the Eighth Circuit affirmed (\href{https://caselaw.findlaw.com/court/us-8th-circuit/114759006.html}{791 F.2d 628 (8th Cir. 1986)}). The Supreme Court granted certiorari to resolve whether Congress could indirectly encourage a uniform national minimum drinking age of 21 through conditional highway grants, even assuming the Twenty-first Amendment barred direct federal regulation of state drinking ages (\href{https://supreme.justia.com/cases/federal/us/483/203/}{483 U.S. at 205--206}).

\section{Issue(s) Presented}

\begin{enumerate}
  \item Whether Congress exceeds its authority under the Spending Clause (U.S. Const. art. I, \S~8, cl. 1) by directing the Secretary of Transportation to withhold 5\% of specified federal highway funds from states that permit persons under age 21 to purchase or publicly possess alcoholic beverages.
  \item Whether Section 2 of the Twenty-first Amendment constitutes an ``independent constitutional bar'' preventing Congress from using conditional spending to induce states to adopt a minimum drinking age that Congress is assumed unable to mandate directly.
  \item Whether conditioning 5\% of federal highway grant funds on raising a state's minimum drinking age to 21 crosses the line from permissible financial encouragement into unconstitutional coercion of state sovereigns.
\end{enumerate}

\section{Holding}

\textbf{No to all issues; judgment affirmed (7--2).} Writing for the Court, Chief Justice Rehnquist held that even assuming Congress lacks power under the Twenty-first Amendment to impose a national minimum drinking age directly (a question the Court expressly declined to decide), 23 U.S.C. \S~158 is a valid exercise of the Spending Clause (\href{https://supreme.justia.com/cases/federal/us/483/203/}{483 U.S. at 206, 212}). Conditioning 5\% of specified federal highway funds on a state's adoption of a 21-year-old minimum drinking age: (1) pursues the ``general welfare''; (2) states the condition unambiguously; (3) is directly related to the federal interest in safe interstate travel; (4) violates no ``independent constitutional bar,'' because the condition does not induce states to engage in unconstitutional conduct; and (5) constitutes ``relatively mild encouragement'' rather than unconstitutional compulsion (\href{https://supreme.justia.com/cases/federal/us/483/203/}{483 U.S. at 207--212}).

\section{Reasoning}

\subsection{Bypassing the Twenty-first Amendment Scope Question and Affirming Hamiltonian Spending Power}

Although both parties focused heavily on whether Section 2 of the Twenty-first Amendment reserves minimum drinking-age regulation exclusively to the states, Chief Justice Rehnquist sidestepped that question entirely (\href{https://supreme.justia.com/cases/federal/us/483/203/}{483 U.S. at 206}; see also \href{https://www.law.cornell.edu/constitution-conan/amendment-21/state-and-federal-regulation-of-minimum-drinking-age}{Cornell LII Constitution Annotated, Amend. XXI}). Relying on \href{https://supreme.justia.com/cases/federal/us/297/1/}{\textit{United States v. Butler}, 297 U.S. 1, 65--66 (1936)}, the Court reaffirmed the broad Hamiltonian view of the Spending Clause: Congress's power to spend public money for public purposes ``is not limited by the direct grants of legislative power found in the Constitution.'' Consequently, Congress may attach conditions to federal funds to attain policy objectives that lie outside Article I's enumerated regulatory fields (\href{https://supreme.justia.com/cases/federal/us/483/203/}{483 U.S. at 206--207}).

\subsection{The Four \textit{Dole} Restrictions Plus the Anti-Coercion Limit}

Synthesizing prior Spending Clause precedent, the Court articulated four general restrictions on conditional grants---followed by a fifth, structural anti-coercion constraint---and found \S~158 satisfied every requirement (\href{https://supreme.justia.com/cases/federal/us/483/203/}{483 U.S. at 207--212}; see also \href{https://www.law.cornell.edu/constitution-conan/article-1/section-8/clause-1/general-welfare-relatedness-and-independent-constitutional-bars}{Cornell LII Constitution Annotated, Art. I, \S~8, cl. 1}):

\begin{enumerate}
  \item \textbf{Pursuit of the ``General Welfare'':} Derived from the text of Article I, \S~8, cl. 1, the expenditure must serve general public purposes. Courts must ``defer substantially to the judgment of Congress'' (\textit{Helvering v. Davis}, 301 U.S. 619, 640, 645 (1937)); indeed, the Court noted in Footnote 2 that it had previously questioned ``whether `general welfare' is a judicially enforceable restriction at all'' (\href{https://supreme.justia.com/cases/federal/us/483/203/}{483 U.S. at 207--208 \& n.2}). Addressing the interstate drunk-driving hazard caused by nonuniform drinking ages easily met this deferential standard.
  \item \textbf{Unambiguous Condition (Clear Notice):} Drawing on \textit{Pennhurst State School and Hospital v. Halderman}, 451 U.S. 1, 17 (1981), Congress must state funding conditions ``unambiguously,'' enabling states ``to exercise their choice knowingly, cognizant of the consequences of their participation.'' Section 158 stated the 21-year-old threshold and the 5\% withholding penalty with total clarity (\href{https://supreme.justia.com/cases/federal/us/483/203/}{483 U.S. at 207--208}).
  \item \textbf{Germaneness / Relatedness:} Conditions might be illegitimate if ``unrelated `to the federal interest in particular national projects or programs''' (\textit{Massachusetts v. United States}, 435 U.S. 444, 461 (1978) (plurality opinion)). Without defining the outer bounds of germaneness (Footnote 3), the Court found \S~158 ``directly related to one of the main purposes for which highway funds are expended---safe interstate travel'' (23 U.S.C. \S~101(b)), because differing state drinking ages incentivized young drivers to cross borders to drink and drive (\href{https://supreme.justia.com/cases/federal/us/483/203/}{483 U.S. at 207--209 \& n.3}).
  \item \textbf{No ``Independent Constitutional Bar'':} This was the centerpiece of South Dakota's challenge. The State argued that the Twenty-first Amendment barred Congress from doing indirectly via spending what it could not do directly. Rehnquist rejected that reading, relying on \href{https://supreme.justia.com/cases/federal/us/297/1/}{\textit{Butler}} and \textit{Oklahoma v. Civil Service Comm'n}, 330 U.S. 127, 143--144 (1947) (upholding Hatch Act funding conditions despite Tenth Amendment objections). The ``independent constitutional bar'' does \textit{not} forbid indirect achievement of objectives beyond Congress's direct regulatory powers; rather, it forbids Congress from using federal funds ``to induce the States to engage in activities that would themselves be unconstitutional''---such as invidious racial discrimination or cruel and unusual punishment. Because South Dakota itself could constitutionally raise its drinking age to 21 without violating anyone's constitutional rights, the Twenty-first Amendment posed no independent bar (\href{https://supreme.justia.com/cases/federal/us/483/203/}{483 U.S. at 209--211}).
  \item \textbf{No Coercion (``Pressure Turns into Compulsion''):} Citing \textit{Steward Machine Co. v. Davis}, 301 U.S. 548, 590 (1937), the Court acknowledged that a financial inducement could theoretically become so severe that ``pressure turns into compulsion.'' Here, however, South Dakota stood to lose only 5\% of funds under specified highway grant programs---a sum later quantified in \href{https://supreme.justia.com/cases/federal/us/567/519/}{\textit{NFIB v. Sebelius}, 567 U.S. 519, 581 (2012)}, as less than half of one percent of South Dakota's total state budget. Rejecting the State's argument that the statute's near-universal compliance proved coercion, Rehnquist emphasized that a conditional grant is not unconstitutional merely because it succeeds; raising the drinking age remained ``the prerogative of the States not merely in theory, but in fact'' (\href{https://supreme.justia.com/cases/federal/us/483/203/}{483 U.S. at 211--212}).
\end{enumerate}

\section{Concurring and Dissenting Opinions}

\subsection{Justice Brennan's Dissent (483 U.S. at 212)}

In a concise, one-paragraph dissent, Justice Brennan agreed with Justice O'Connor that setting the minimum age for purchasing liquor falls squarely within the powers reserved to the states by Section 2 of the Twenty-first Amendment. Unlike O'Connor, however, Brennan treated the Twenty-first Amendment itself as a substantive limit on conditional spending: because the Amendment vests states with that specific constitutional power, ``Congress cannot condition a federal grant in a manner that abridges this right. The Amendment, itself, strikes the proper balance between federal and state authority'' (\href{https://supreme.justia.com/cases/federal/us/483/203/}{483 U.S. at 212}).

\subsection{Justice O'Connor's Dissent (483 U.S. at 212--218)}

Justice O'Connor accepted the majority's four-part framework---conceding that \S~158 satisfied the general welfare and clear-notice requirements and assuming \textit{arguendo} that the Twenty-first Amendment was not an independent constitutional bar---but sharply dissented on the \textbf{relatedness (germaneness)} prong (\href{https://supreme.justia.com/cases/federal/us/483/203/}{483 U.S. at 212--214}).

\begin{itemize}
  \item \textbf{Over- and Under-Inclusiveness:} O'Connor argued that conditioning highway construction funds on a uniform drinking age of 21 was ``far too over- and under-inclusive'' to be reasonably related to highway safety: overinclusive because it bars teenagers from drinking even when they are not driving on interstate highways, and underinclusive because, as congressional floor debates showed, 84\% of alcohol-related highway fatalities involved drivers over age 21 (\href{https://supreme.justia.com/cases/federal/us/483/203/}{483 U.S. at 214--215}).
  \item \textbf{The ``Condition vs. Regulation'' Bright Line:} Adopting the test urged by the National Conference of State Legislatures as \textit{amicus curiae} and grounded in \href{https://supreme.justia.com/cases/federal/us/297/1/}{\textit{United States v. Butler}, 297 U.S. at 73, 78}, O'Connor proposed a structural boundary: Congress may specify \textit{how the federal money itself should be spent} (a legitimate condition), but it may not use the Spending Clause to impose behavioral mandates outside the expenditure's administration (a regulation valid only under an enumerated regulatory power). Because \S~158 did not govern how highway dollars were spent, it was an unconstitutional regulation of liquor sales barred under the Commerce Clause by Section 2 of the Twenty-first Amendment (\href{https://supreme.justia.com/cases/federal/us/483/203/}{483 U.S. at 215--218}).
  \item \textbf{Slippery-Slope Federalism Warning:} O'Connor warned that under the majority's loose relatedness inquiry, Congress could condition highway money on moving a state capital or regulating almost any aspect of a state's social, political, or economic life, eroding the fundamental principle of \textit{McCulloch v. Maryland}, 17 U.S. (4 Wheat.) 316, 405 (1819), that the federal government remains one of enumerated powers (\href{https://supreme.justia.com/cases/federal/us/483/203/}{483 U.S. at 215, 218}).
\end{itemize}

\section{Key Quotations}

\begin{enumerate}
  \item \textbf{Spending Power Extends Beyond Enumerated Regulatory Powers (Rehnquist, C.J., \href{https://supreme.justia.com/cases/federal/us/483/203/}{483 U.S. at 207}):}
  \begin{quote}
    ``Thus, objectives not thought to be within Article I's `enumerated legislative fields,' \textit{id.} at 297 U. S. 65, may nevertheless be attained through the use of the spending power and the conditional grant of federal funds.''
  \end{quote}

  \item \textbf{Meaning of the ``Independent Constitutional Bar'' (Rehnquist, C.J., \href{https://supreme.justia.com/cases/federal/us/483/203/}{483 U.S. at 210}):}
  \begin{quote}
    ``These cases establish that the `independent constitutional bar' limitation on the spending power is not, as petitioner suggests, a prohibition on the indirect achievement of objectives which Congress is not empowered to achieve directly. Instead, we think that the language in our earlier opinions stands for the unexceptionable proposition that the power may not be used to induce the States to engage in activities that would themselves be unconstitutional.''
  \end{quote}

  \item \textbf{The Anti-Coercion Boundary and ``Relatively Mild Encouragement'' (Rehnquist, C.J., \href{https://supreme.justia.com/cases/federal/us/483/203/}{483 U.S. at 211--212}):}
  \begin{quote}
    ``Our decisions have recognized that, in some circumstances, the financial inducement offered by Congress might be so coercive as to pass the point at which `pressure turns into compulsion.' . . . Here Congress has offered relatively mild encouragement to the States to enact higher minimum drinking ages than they would otherwise choose. But the enactment of such laws remains the prerogative of the States not merely in theory, but in fact.''
  \end{quote}

  \item \textbf{Justice O'Connor on the Limit of Relatedness---Condition vs. Regulation (O'Connor, J., dissenting, \href{https://supreme.justia.com/cases/federal/us/483/203/}{483 U.S. at 215, 218}):}
  \begin{quote}
    ``When Congress appropriates money to build a highway, it is entitled to insist that the highway be a safe one. But it is not entitled to insist as a condition of the use of highway funds that the State impose or change regulations in other areas of the State's social and economic life because of an attenuated or tangential relationship to highway use or safety. . . . Rather than a condition determining how federal highway money shall be expended, it is a regulation determining who shall be able to drink liquor. As such, it is not justified by the spending power.''
  \end{quote}

  \item \textbf{Justice Brennan on Twenty-first Amendment Sovereignty (Brennan, J., dissenting, \href{https://supreme.justia.com/cases/federal/us/483/203/}{483 U.S. at 212}):}
  \begin{quote}
    ``Since States possess this constitutional power, Congress cannot condition a federal grant in a manner that abridges this right. The Amendment, itself, strikes the proper balance between federal and state authority.''
  \end{quote}
\end{enumerate}

\section{Doctrinal Significance and Modern Legacy (Through 2026)}

For nearly four decades, \textit{South Dakota v. Dole} has served as the canonical framework governing conditional federal spending---the primary mechanism through which Congress shapes state policy in areas such as education, healthcare, transportation, and civil rights without violating the anti-commandeering doctrine of \textit{New York v. United States}, 505 U.S. 144 (1992), and \textit{Printz v. United States}, 521 U.S. 898 (1997). Yet while \textit{Dole} initially appeared to give Congress an almost unchecked ``end run'' around federalism limits (especially after \textit{United States v. Lopez}, 514 U.S. 549 (1995)), Supreme Court jurisprudence from 2012 through 2026 has dramatically tightened three distinct edges of the \textit{Dole} doctrine:

\begin{itemize}
  \item \textbf{Anti-Coercion Given Real Teeth (\textit{NFIB v. Sebelius}, 2012):} For twenty-five years after \textit{Dole}, no federal statute was struck down as coercive. That changed in \href{https://supreme.justia.com/cases/federal/us/567/519/}{\textit{National Federation of Independent Business v. Sebelius}, 567 U.S. 519 (2012)}, where seven Justices (Chief Justice Roberts's controlling plurality joined by Justices Breyer and Kagan, plus the four-Justice joint dissent) held that the Affordable Care Act's Medicaid expansion crossed \textit{Dole}'s coercion line. Contrasting South Dakota's 5\% highway-fund cut (``less than half of one percent of South Dakota's budget'') with the ACA's threat to strip \textit{all} existing Medicaid funding (over 10\% of an average state's total budget), Roberts famously labeled the Medicaid condition ``a gun to the head'' rather than ``relatively mild encouragement'' (\href{https://www.law.cornell.edu/constitution-conan/article-1/section-8/clause-1/anti-coercion-requirement-and-spending-clause}{Cornell LII Constitution Annotated, Anti-Coercion Requirement}).
  \item \textbf{Independent Constitutional Bars and Direct Protection of Federal Funds:} In \href{https://supreme.justia.com/cases/federal/us/539/194/}{\textit{United States v. American Library Ass'n, Inc.}, 539 U.S. 194 (2003)}, the Court upheld library Internet-filtering conditions under \textit{Dole}'s independent-bar framework, whereas \href{https://www.law.cornell.edu/supct/pdf/03-44P.ZS}{\textit{Sabri v. United States}, 541 U.S. 600 (2004)}, distinguished \textit{Dole} to uphold a federal anti-bribery statute (18 U.S.C. \S~666(a)(2)) under the Spending and Necessary and Proper Clauses as a direct safeguard against graft in federally funded entities.
  \item \textbf{The ``Contract Analogy'' Tightens Clear Notice, Remedies, and Consent (2022--2026):} In the 2020s, the Roberts Court transformed \textit{Dole}'s second prong (clear notice under \textit{Pennhurst}) into a rigorous contract-law constraint on who can be sued and what remedies apply under Spending Clause statutes:
  \begin{itemize}
    \item In \href{https://www.supremecourt.gov/opinions/21pdf/20-219_1b82.pdf}{\textit{Cummings v. Premier Rehab Keller, P.L.L.C.}, 596 U.S. 212 (2022)}, the Court held that emotional distress damages are unavailable in private suits under Spending Clause antidiscrimination laws (the Rehabilitation Act and ACA) because funding recipients lack clear notice of remedies not ``traditionally available in suits for breach of contract.''
    \item In \href{https://www.supremecourt.gov/opinions/22pdf/599us1r36_feah.pdf}{\textit{Health and Hospital Corp. of Marion County v. Talevski}, 599 U.S. 166 (2023)}, the Court reaffirmed that Spending Clause statutes \textit{can} create privately enforceable rights under 42 U.S.C. \S~1983 (upholding FNHRA nursing-home rights), but emphasized a ``demanding bar'' requiring unambiguous rights-creating text.
    \item Two years later, in \href{https://supreme.justia.com/cases/federal/us/606/23-1275/}{\textit{Medina v. Planned Parenthood South Atlantic}, 606 U.S. 357 (2025)}, the Court narrowed \S~1983 enforcement sharply, holding that Medicaid's ``any-qualified-provider'' provision lacks the explicit rights-creating language required to overcome the presumption that funding termination---not private litigation---is the typical remedy for noncompliance with spending legislation.
    \item Most recently, in \href{https://www.supremecourt.gov/opinions/25pdf/609us1r55_h315.pdf}{\textit{Landor v. Louisiana Department of Corrections and Public Safety}, 609 U.S. 59 (2026)}, the Court held that neither the Spending Clause nor the Necessary and Proper Clause authorizes RLUIPA personal-capacity money-damages suits against individual state prison officials. Citing \textit{Dole}, \textit{Cummings}, and \textit{Medina}, the Court stressed that because the Spending Clause confers no authority to ``regulate directly'' (\textit{Dole}, 483 U.S. at 209), spending conditions bind only parties who ``voluntarily and knowingly'' consent---a state agency's acceptance of federal funds cannot bind nonconsenting individual employees (\href{https://www.everycrsreport.com/files/2026-10-07_R46827_c76ab672b14066efe09e7303065f8022faea080c.pdf}{CRS Report R46827, Oct. 7, 2026}).
  \end{itemize}
  \item \textbf{Executive-Imposed Grant Conditions (2025--2026):} Beyond congressional statutes, federal litigation through 2026 over executive-branch funding conditions (e.g., conditioning local and municipal grants on immigration or policy compliance, as in \href{https://www.epa.gov/system/files/documents/2026-10/port-of-seattle-et-al-v-blanche-et-al-complaint.pdf}{\textit{Port of Seattle v. Blanche} (W.D. Wash. filed Oct. 2026)}) routinely invokes \textit{Dole}'s relatedness, clear-notice, and anti-coercion prongs alongside separation-of-powers limits, underscoring that agencies cannot impose conditions Congress itself neither authorized nor could constitutionally attach (\href{https://www.everycrsreport.com/files/2026-10-07_R46827_c76ab672b14066efe09e7303065f8022faea080c.pdf}{CRS Report R46827, Oct. 7, 2026}).
\end{itemize}

\section{Pedagogical Critiques and Exam Hypotheticals}

\subsection{Key Pedagogical Critiques}

\begin{enumerate}
  \item \textbf{The ``Trojan Horse'' Problem for Enumerated Powers:} Scholars and jurists have long noted an irony: Chief Justice Rehnquist---the principal architect of the 1990s federalism revival in \textit{Lopez} and \textit{Morrison}---authored the very opinion in \textit{Dole} that allows Congress to circumvent Commerce Clause limits by ``buying'' state regulatory compliance. Because federal taxation extracts revenue from the same taxpayers who fund state treasuries, states face intense political pressure to accept ``free'' federal money even when they oppose the attached policy.
  \item \textbf{Indeterminacy of Relatedness and Coercion:} Neither \textit{Dole} nor \textit{NFIB} supplies a bright-line metric for when a condition becomes too tangential (relatedness) or too large (coercion). If 5\% of highway funds (less than 0.5\% of a state's budget) is ``mild encouragement'' and 100\% of Medicaid funds (over 10\% of a state's budget) is a ``gun to the head,'' what about 25\% of federal education funds or 40\% of highway funds? Likewise, \textit{Dole} Footnote 3 expressly left the outer bounds of germaneness undefined, leaving lower courts to wrestle with Justice O'Connor's ``condition vs. regulation'' critique.
  \item \textbf{Asymmetry of the ``Independent Constitutional Bar'':} Rehnquist's reading protects individual-rights provisions (First, Eighth, and Fourteenth Amendments) from indirect subversion, because a state itself cannot violate them. Yet it strips structural federalism provisions (like the Tenth Amendment and, as assumed in \textit{Dole}, the Twenty-first Amendment) of independent force as bars, treating state sovereignty as fully waivable for cash.
\end{enumerate}

\subsection{Exam Hypotheticals (Including ``Opposite Side'' Contrast)}

\subsubsection{Hypothetical 1 (Classic \textit{Dole} / \textit{NFIB} Synthesis---Relatedness and Coercion)}
Congress enacts the \textit{Safe Campuses Act}, directing the Secretary of Education to withhold 100\% of a state's existing federal K--12 Title I funding (8\% of the average state's total operating budget) unless the state enacts a criminal ban on all firearm possession within 1,000 feet of any public university campus and directs state police to enforce it.

\textbf{Analytical Roadmap:}
\begin{itemize}
  \item \textit{General Welfare \& Clear Notice:} Satisfied; public safety receives near-total deference under \href{https://supreme.justia.com/cases/federal/us/483/203/}{\textit{Dole}}, and the condition is unambiguous.
  \item \textit{Relatedness:} Arguably satisfied under Chief Justice Rehnquist's deferential standard (campus safety relates to education funding), but vulnerable under Justice O'Connor's dissent because it regulates criminal law around \textit{universities} using \textit{K--12} funding rather than specifying how Title I funds are spent.
  \item \textit{Independent Constitutional Bar:} If a blanket 1,000-foot buffer zone violates the Second Amendment as applied to public sidewalks outside sensitive places, the condition fails \textit{Dole}'s fourth prong regardless of funding size.
  \item \textit{Coercion:} Unconstitutional under \href{https://supreme.justia.com/cases/federal/us/567/519/}{\textit{NFIB v. Sebelius}, 567 U.S. 519 (2012)}. Threatening 100\% of an existing funding stream totaling 8\% of a state's budget is a ``gun to the head,'' far beyond \textit{Dole}'s 5\% highway cut (less than 0.5\% of South Dakota's budget).
\end{itemize}

\subsubsection{Hypothetical 2 (``Opposite Side'' Contrast---Valid Conditional Grant vs. Impermissible Spending Remedy)}
Compare two federal statutes promoting electric-vehicle (EV) charging infrastructure:
\begin{itemize}
  \item \textbf{Statute A (Valid Conditional Grant---The \textit{Dole} Side):} Congress creates a new \$10 billion Federal EV Corridor Grant Program. Participating states must (1) spend the funds exclusively on public fast-charging stations along interstate highways and (2) require EV-ready wiring in newly constructed commercial parking garages. Nonparticipating states forgo only the new grant and lose zero existing highway funds.
  \item \textbf{Statute B (Invalid Spending Condition / Remedy---The ``Opposite Side''):} Congress amends an existing \$50 billion Highway Trust Fund program to (1) authorize motorists to sue individual state highway engineers in their \textit{personal capacities} for emotional distress damages whenever a state EV charger is out of service for over 48 hours, and (2) withhold 80\% of existing federal highway funds from any state that declines to ban gasoline-powered lawnmowers statewide.
\end{itemize}

\textbf{Contrast Analysis:}
\begin{itemize}
  \item \textbf{Why Statute A Survives:} It pursues the general welfare, gives clear prospective notice, relates directly to interstate EV transportation (valid under \textit{Dole}'s majority even though requirement (2) regulates building codes), induces no unconstitutional conduct, and avoids coercion by risking no pre-existing funds.
  \item \textbf{Why Statute B Fails:} Cutting 80\% of existing highway funds to force a lawnmower ban fails both \textbf{coercion} (\href{https://supreme.justia.com/cases/federal/us/567/519/}{\textit{NFIB}}) and \textbf{relatedness} (\href{https://supreme.justia.com/cases/federal/us/483/203/}{\textit{Dole}} prong 3). Moreover, \textbf{emotional distress damages} violate the contract-law notice rule of \href{https://www.supremecourt.gov/opinions/21pdf/20-219_1b82.pdf}{\textit{Cummings v. Premier Rehab Keller, P.L.L.C.}, 596 U.S. 212 (2022)}, and \textbf{personal-capacity damages against nonconsenting state employees} violate \href{https://www.supremecourt.gov/opinions/25pdf/609us1r55_h315.pdf}{\textit{Landor v. Louisiana Department of Corrections and Public Safety}, 609 U.S. 59 (2026)}, because a state's receipt of funds cannot bind individual officers who did not consent to the spending contract.
\end{itemize}

\end{document}
